% SPDX-FileCopyrightText: 2026 Christoph Maier
% SPDX-License-Identifier: Apache-2.0
%
% Running log of the class-AB pad driver revision (LaTeX version of log.md).
% Body converted from log.md with pandoc 3.1.3 (gfm -> latex, --shift-heading-level-by=-1,
% --columns=50), tables and '~' post-processed; preamble written by hand in the style of ../../../esd_protection/*.tex.
% Build: pdflatex log.tex (twice). Packages: standard LaTeX tools, booktabs, longtable, calc,
% hyperref, xurl, microtype, textcomp, geometry. No IEEEtran, no siunitx.
\documentclass[11pt,a4paper]{article}

\usepackage[utf8]{inputenc}
\usepackage[a4paper,margin=22mm]{geometry}
\usepackage[expansion=false]{microtype}
\usepackage{amsmath,amssymb}
\usepackage{textcomp}
\usepackage{longtable,booktabs,array,calc}
\usepackage[hidelinks]{hyperref}
\usepackage{xurl}

% characters outside T1 used in log.md
\DeclareUnicodeCharacter{03A9}{\ensuremath{\Omega}}
\DeclareUnicodeCharacter{03BA}{\ensuremath{\kappa}}
\DeclareUnicodeCharacter{03C3}{\ensuremath{\sigma}}
\DeclareUnicodeCharacter{2192}{\ensuremath{\rightarrow}}
\DeclareUnicodeCharacter{2212}{\ensuremath{-}}
\DeclareUnicodeCharacter{2213}{\ensuremath{\mp}}
\DeclareUnicodeCharacter{221A}{\ensuremath{\surd}}
\DeclareUnicodeCharacter{2225}{\ensuremath{\parallel}}
\DeclareUnicodeCharacter{2248}{\ensuremath{\approx}}
\DeclareUnicodeCharacter{2264}{\ensuremath{\le}}
\DeclareUnicodeCharacter{2265}{\ensuremath{\ge}}

% what pandoc's body expects
\providecommand{\tightlist}{\setlength{\itemsep}{0pt}\setlength{\parskip}{0pt}}
\setlength{\emergencystretch}{3em}
\setlength{\parindent}{0pt}
\setlength{\parskip}{4pt plus 2pt}
\renewcommand{\arraystretch}{1.1}
\setcounter{secnumdepth}{0}
% allow line breaks after underscores (file and net names in \texttt)
\renewcommand{\_}{\textunderscore\hspace{0pt}}

\title{Running log --- improvements to the class-AB pad driver\\[2pt]
  \large \texttt{sg13cmos5l\_cm\_ip\_\_single2diff2single}, \texttt{design\_considerations/class\_ab\_pad\_driver/improvements}}
\author{}
\date{2 October 2026}

\begin{document}
\maketitle
\small
Session 2026-10-02, Claude (configured model
\texttt{claude-opus-5-5}; the serving model may
differ). Newest entries at the bottom. Times are
Europe/Berlin, taken from the session transcript
(corrected 15:01: the first versions of several
headings carried guessed times, up to 1¼ h too
late).

Task as given: area, current and loop gain of
\texttt{d2s\_miller} / \texttt{d2s\_loadcomp} are
unsatisfactory, probably because of the linear
passives (rhigh degeneration, MOM Miller caps).
Look in Edwards, Minch, Cauwenberghs, Hasler
(DoNotLitter/TimEdwards, LogDomain, FloatingGate)
and Sarpeshkar 1997 for (1) alternatives to
resistor source degeneration in the DDA, (2)
compensation without linear BEOL capacitance, (3)
cascoding with the output devices either doubling
as ESD clamps or as separate devices in the slot,
(4) externally-linear / internally-nonlinear
circuits whose internal nonlinearity is not the
subthreshold exponential. End result: proposed
improvements to \texttt{class\_ab\_pad\_driver}.

Nothing is written to
\texttt{\textasciitilde{}/DoNotLitter}. Papers are
cited, not copied (chatlog/ref rule).

\hypertarget{1410--baseline-taken-from-the-existing-note-and-the-first-cace-run}{%
\section{14:10 --- baseline taken from the
existing note and the first CACE
run}\label{1410--baseline-taken-from-the-existing-note-and-the-first-cace-run}}

What the CACE run
(\texttt{verification/cace/first\_run\_summary.md})
adds to the design note:

\begin{itemize}
\tightlist
\item
  Gain spread 0.476\ldots0.554 over PVT + rhigh
  corners, 0.41\ldots0.61 in process MC. Cause:
  the gain is (R2 + 2/gm)/(R1 + 2/gm); R tracks R,
  but gm does not track R.
\item
  Output offset in mismatch MC −32\ldots+49 mV (4
  degenerated pairs; offset referred through 1/Gm
  of degenerated pairs is large).
\item
  INL 469 mV worst corner; THD −15 dB at 1.25 V
  output (pairs run out of I\_tail·R range).
\item
  Output noise 1.6 µV/√Hz at 1 kHz (degeneration
  resistors + low Gm).
\item
  Area: the two MOM Miller caps (2 × 31 × 31 µm)
  dominate the non-frame area (1920 of
  $\sim$2500 µm²).
\item
  Current: 427 µA, of which 240 µA output-stage
  quiescent, $\sim$180 µA front end +
  bias.
\end{itemize}

\hypertarget{1425--literature-read-so-far-details-in-literature_notesmd}{%
\section{\texorpdfstring{14:25 --- literature read
so far (details in
\texttt{literature\_notes.md})}{14:25 --- literature read so far (details in literature\_notes.md)}}\label{1425--literature-read-so-far-details-in-literature_notesmd}}

\begin{itemize}
\tightlist
\item
  Sarpeshkar 1997 thesis (scanned,
  OCR\textquotesingle d locally on the linked
  computer), ch. 3 = the wide-linear-range OTA:
  well input, diode source degeneration, gate
  degeneration through the mirror diode, bump
  linearization (w = 2 cancels the cubic term). No
  resistors.
\item
  Edwards thesis §3.5--3.10, Edwards \&
  Cauwenberghs 2000: ELIN/log-domain; names the
  square law of strong inversion as the other
  usable internal nonlinearity; matching = "same
  surround".
\item
  Minch thesis ch. 7: a source follower is linear
  whatever the device law (MITE via FG source
  follower, valid above threshold).
\item
  Hasler thesis ch. 4--5, Sánchez-Sinencio FG
  slides: capacitive dividers / FG inputs widen
  linear range, but need linear caps and charge
  control (tunneling/injection or UV).
\end{itemize}

\hypertarget{1425--simulation-environment}{%
\section{14:25 --- simulation
environment}\label{1425--simulation-environment}}

Cloud container: ngspice-42 (Ubuntu), IHP-Open-PDK
\texttt{dev} sparse checkout, OpenVAF-reloaded
built from source with LLVM 18 for psp103 /
r3\_cmc / cap\_cmomi. (in progress)

\hypertarget{1428--simulation-environment-ready-baseline-reproduced}{%
\section{14:28 --- simulation environment ready,
baseline
reproduced}\label{1428--simulation-environment-ready-baseline-reproduced}}

ngspice-42 + IHP-Open-PDK \texttt{dev} models +
psp103 / r3\_cmc / cap\_cmomi compiled with
OpenVAF-reloaded (LLVM 18, OSDI minor version
patched 4 → 3 for ngspice 42, as in the first
session). \texttt{run\_d2s.py} functions reproduce
\texttt{results.txt} exactly (A, 1 kΩ ∥ 100 pF:
gain 0.5005, 0.18 mV, nonlin 5.52 mV, T0 66.5 dB,
fc 1.62 MHz, PM 74.8°, THD 10 kHz −51.1 dB).

\hypertarget{1440--dda-unit-characterization-simunitsspice-test-bench-now-in-run_improvementspy-units}{%
\section{\texorpdfstring{14:40 --- DDA unit
characterization (\texttt{sim/units.spice}; test
bench now in
\texttt{run\_improvements.py\ units})}{14:40 --- DDA unit characterization (sim/units.spice; test bench now in run\_improvements.py units)}}\label{1440--dda-unit-characterization-simunitsspice-test-bench-now-in-run_improvementspy-units}}

Each unit is one split-tail PMOS pair; x = gp − gn
swept ±0.8 V at input CM 1.40 / 1.65 / 1.90 V.

\begin{longtable}[]{@{}>{\raggedright\arraybackslash}p{(\linewidth - 10\tabcolsep) * \real{0.1053}}>{\raggedright\arraybackslash}p{(\linewidth - 10\tabcolsep) * \real{0.3070}}>{\raggedright\arraybackslash}p{(\linewidth - 10\tabcolsep) * \real{0.2234}}>{\raggedright\arraybackslash}p{(\linewidth - 10\tabcolsep) * \real{0.1720}}>{\raggedright\arraybackslash}p{(\linewidth - 10\tabcolsep) * \real{0.1923}}@{}}
\toprule\noalign{}
unit & degeneration & Gm0 @ CM 1.40 / 1.65 / 1.90
& f(0.5)/(Gm0·0.5) & CM sensitivity of Gm \\
\midrule\noalign{}
\endhead
\bottomrule\noalign{}
\endlastfoot
unit\_r & rhigh 0.5 × 36.4 µm (108 kΩ), I\_t = 2 ×
10 µA & 15.54 / 15.55 / 15.58 µS & 0.996 & 0.5
\%/V \\
unit\_t & hv PMOS 0.5/3 in triode, gate at vss &
12.43 / 15.54 / 18.40 µS & 0.992 & 77 \%/V \\
unit\_w & Sarpeshkar well input, 4/4, gates at
vss, 1 × 20 µA & 11.82 / 11.02 / 10.35 µS & 0.991
& −27 \%/V \\
\end{longtable}

The MOS alternatives are linear enough on their
own; their problem is common-mode sensitivity
(triode R depends on V\_SG = source − gate;
well-input κ depends on well-to-gate voltage, as
Sarpeshkar\textquotesingle s §3.3.1 reports). Why
that matters for the next step: see 14:43 (a
leftover "15:05" from the guessed timestamps,
corrected 16:15).

\hypertarget{1443--matched-pair-dda-mp-dda}{%
\section{14:43 --- matched-pair DDA
(MP-DDA)}\label{1443--matched-pair-dda-mp-dda}}

Idea: four identical units, A (vinp, vref), B
(vref, vinn), C1 = C2 (vref, vfb): f(vinp − vref)
+ f(vref − vinn) = 2 f(vfb − vref). With matched
units the solution is vout − vref = (vinp −
vinn)/2 for \emph{any} odd monotonic f (exact for
input CM = vref; for a CM error c the error is
second order in c). Gain ½ = unit count ratio; no
R1/R2, no gm term. This is static ELIN: the input
units compress, the feedback units invert the same
law. Caveat found analytically: A and C sit at CM
vref + x/2, B at vref − x/2. A unit whose Gm
depends on CM (k per volt) gives v ≈ u − k u²/2 →
unit\_t ≈ 96 mV, unit\_w ≈ 34 mV of second-order
error at u = 0.5 V, unit\_r ≈ 0.6 mV. So in this
topology the poly resistor stays the right
degenerator. Its area is $\sim$18 µm² per
unit, it never was the area problem.

First simulation, everything else unchanged (fold,
mirror, class-AB, frames, MOM caps, bias):

\begin{longtable}[]{@{}>{\raggedright\arraybackslash}p{(\linewidth - 12\tabcolsep) * \real{0.2697}}>{\raggedright\arraybackslash}p{(\linewidth - 12\tabcolsep) * \real{0.1364}}>{\raggedright\arraybackslash}p{(\linewidth - 12\tabcolsep) * \real{0.1008}}>{\raggedright\arraybackslash}p{(\linewidth - 12\tabcolsep) * \real{0.1088}}>{\raggedright\arraybackslash}p{(\linewidth - 12\tabcolsep) * \real{0.2097}}>{\raggedright\arraybackslash}p{(\linewidth - 12\tabcolsep) * \real{0.1745}}@{}}
\toprule\noalign{}
variant & load & gain & nonlin & T0 / fc / PM &
THD 10 k / 100 kHz \\
\midrule\noalign{}
\endhead
\bottomrule\noalign{}
\endlastfoot
baseline A & 1 k ∥ 100 p & 0.5005 & 5.52 mV & 66.5
dB / 1.62 MHz / 74.8° & −51.1 / −50.2 dB \\
MP-DDA, MOM & 1 k ∥ 100 p & 0.4998 & 1.04 mV &
66.5 dB / 1.62 MHz / 74.7° & −60.2 / −55.5 dB \\
MP-DDA, MOM & 100 p & 0.5000 & 1.02 mV & 96.5 dB /
2.03 MHz / 66.2° & −60.0 / −59.8 dB \\
MP-DDA, hv-PMOS accumulation cap 14 × 14 µm & 100
p & 0.5000 & 1.02 mV & 96.7 dB / 2.80 MHz / 53.4°
& −60.0 / −59.9 dB \\
\end{longtable}

Same current (427 µA) and tail budget (8 × 10 µA =
baseline 4 × 20 µA). 9 dB better THD, 5× lower DC
nonlinearity. The MOS cap (196 µm² vs 961 µm² MOM)
does not degrade THD despite a 2× C--V variation
over 0.15\ldots1.25 V (0.37 → 0.75 pF); it is just
smaller than 1 pF on average, so it needs resizing
for the PM.

\hypertarget{1448--mos-capacitor-sizing-corners-mismatch}{%
\section{14:48 --- MOS capacitor sizing, corners,
mismatch}\label{1448--mos-capacitor-sizing-corners-mismatch}}

\textbf{MOS accumulation cap size} (MP-DDA, f\_c /
PM): 18 × 18 µm matches the 31 × 31 µm MOM pair
almost exactly (100 pF: 2.00 MHz / 67.0° vs 2.03
MHz / 66.2°; 300 pF: 46.3° vs 45.5°; 1 nF: 26.9°
vs 26.4°). Area per side 324 µm² instead of 961
µm² (3.0×), no metal used.

\textbf{Corners} (dc gain fitted over
\textbar vd\textbar{} ≤ 0.2 V, 1 kΩ ∥ 100 pF;
nonlin = max deviation over ±1 V):

\begin{longtable}[]{@{}>{\raggedright\arraybackslash}p{(\linewidth - 6\tabcolsep) * \real{0.2725}}>{\raggedright\arraybackslash}p{(\linewidth - 6\tabcolsep) * \real{0.3669}}>{\raggedright\arraybackslash}p{(\linewidth - 6\tabcolsep) * \real{0.3605}}@{}}
\toprule\noalign{}
corner & baseline A gain / nonlin & MP-DDA gain /
nonlin \\
\midrule\noalign{}
\endhead
\bottomrule\noalign{}
\endlastfoot
tt 27 & 0.5005 / 5.52 mV & 0.4998 / 1.04 mV \\
ss 125 (res wcs) & 0.5119 / 8.07 & 0.4988 /
1.56 \\
ff −40 (res bcs) & 0.4968 / 3.30 & 0.4998 /
0.66 \\
sf / fs 27 & 0.4998 / 4.11, 0.5013 / 8.07 & 0.4998
/ 0.82, 0.4997 / 1.37 \\
ss 27 / ff 27 & 0.5022 / 12.69, 0.4990 / 3.30 &
0.4998 / 1.91, 0.4997 / 0.69 \\
tt res wcs / bcs & 0.4895 / 7.09, 0.5148 / 5.07 &
0.4997 / 1.48, 0.4998 / 0.66 \\
tt 125 / tt −40 & 0.5249 / 5.18, 0.4870 / 7.15 &
0.4983 / 1.10, 0.4998 / 1.53 \\
\end{longtable}

Systematic gain spread −2.6\ldots+5.0 \% →
−0.34\ldots0 \%. Loop gain / PM unchanged.
(Corrected 15:25: this line first said "±3.8 \%",
which was half the spread, not its range.)

\textbf{Mismatch MC} (30 seeds,
\texttt{mos\_tt\_mismatch} +
\texttt{res\_typ\_mismatch}, mm\_ok=1 on all MOS
and rhigh): baseline offset σ 14.4 mV, MP-DDA σ
14.5 mV --- the front-end change does not touch
the offset. Budget by device group (mismatch on
one group at a time, 20 seeds): folding sinks 11.5
mV, tails 4.2, mirror 2.3, pairs 1.7, everything
else \textless{} 0.4. Same mechanism Sarpeshkar
describes (§3.3.4: current mismatch /
transconductance; low Gm/I front ends are
offset-prone, and the mirrors, not the input
devices, matter most).

\hypertarget{1453--d2s_mp2-same-topology-resized-for-current-and-offset}{%
\section{\texorpdfstring{14:53 ---
\texttt{d2s\_mp2}: same topology, resized for
current and
offset}{14:53 --- d2s\_mp2: same topology, resized for current and offset}}\label{1453--d2s_mp2-same-topology-resized-for-current-and-offset}}

\begin{itemize}
\tightlist
\item
  units at 2 × 5 µA (long-L tails 5/6), R = 150 kΩ
  (rhigh 0.5 × 50.6 µm), I\_t·R = 0.75 V
\item
  folding sinks 30 µA, 36/6 (lower gm/I), mirror
  24/4 at 10 µA
\item
  bias fixture \texttt{d2s\_bias\_lp}: reference
  currents 5 + 5 + 2 + 2 + 5 + 5 µA instead of 20
  + 50 + 5 + 5 + 5 + 5
\item
  headroom limit found on the way: a ≈ 2.6 V and b
  ≈ 0.7 V are the output-device gates, so the
  mirror and the sinks each get only
  $\sim$0.35 V; their overdrive
  can\textquotesingle t be raised for matching,
  only their area.
\item
  all devices ≥ 70 mV from saturation edge at vd =
  0 (\texttt{oppoint.py})
\end{itemize}

\begin{longtable}[]{@{}>{\raggedright\arraybackslash}p{(\linewidth - 6\tabcolsep) * \real{0.3258}}>{\raggedright\arraybackslash}p{(\linewidth - 6\tabcolsep) * \real{0.2964}}>{\raggedright\arraybackslash}p{(\linewidth - 6\tabcolsep) * \real{0.3778}}@{}}
\toprule\noalign{}
& baseline A & d2s\_mp2, MOS 16 × 16 \\
\midrule\noalign{}
\endhead
\bottomrule\noalign{}
\endlastfoot
I\_DD (incl. bias fixture) & 427 µA & 297 µA \\
I\_Q P/N & 240 µA & 213 µA \\
compensation area & 2 × 961 µm² MOM (M1--M4) & 2 ×
256 µm² hv-PMOS \\
T0 1 k ∥ 100 p / 100 p / 50 Ω & 66.5 / 96.5 / 40.7
dB & 69.7 / 99.6 / 43.9 dB \\
fc / PM, 100 p & 2.02 MHz / 66° & 1.76 MHz /
63° \\
THD 10 kHz (1 k ∥ 100 p) & −51.1 dB & −58.7 dB \\
THD 100 kHz (1 k ∥ 100 p) & −50.2 dB & −51.4 dB \\
nonlin (±1 V vd) & 5.52 mV & 1.24 mV \\
offset σ (mismatch, 30) & 14.4 mV & 5.3 mV \\
gain σ (mismatch, 30) & 1.6 \% & 1.1 \% (0.84 \%
from rhigh unit mismatch) \\
\end{longtable}

Gain σ is now set by rhigh area (0.5 × 50.6 µm per
unit): 4× the area halves it.

\hypertarget{1457--the-right-figure-of-merit-for-a-dda-unit-simrun_improvementspy-units}{%
\section{\texorpdfstring{14:57 --- the right
figure of merit for a DDA unit
(\texttt{sim/run\_improvements.py\ units})}{14:57 --- the right figure of merit for a DDA unit (sim/run\_improvements.py units)}}\label{1457--the-right-figure-of-merit-for-a-dda-unit-simrun_improvementspy-units}}

The first unit test swept both inputs around a
common CM. In the MP-DDA one input of every unit
sits at vref, so the relevant comparison is fA(x)
= I(vref + x, vref) against fB(x) = I(vref, vref −
x): the same differential voltage, but the
pair\textquotesingle s sources sit x lower in B.
Solving fA(u) + fB(u) = 2 fA(v) from the unit
curves alone predicts the whole
driver\textquotesingle s static error; for unit\_r
it gives −1.04 / −1.02 mV at u = ∓0.5 V, which is
the full-circuit 1.04 mV.

\begin{longtable}[]{@{}>{\raggedright\arraybackslash}p{(\linewidth - 6\tabcolsep) * \real{0.3839}}>{\raggedright\arraybackslash}p{(\linewidth - 6\tabcolsep) * \real{0.2836}}>{\raggedright\arraybackslash}p{(\linewidth - 6\tabcolsep) * \real{0.3325}}@{}}
\toprule\noalign{}
unit & own compression at 0.5 V & MP-DDA error at
u = −0.5 / +0.5 V \\
\midrule\noalign{}
\endhead
\bottomrule\noalign{}
\endlastfoot
unit\_r (split tails + rhigh) & 0.5 \% & −1.0 /
−1.0 mV \\
unit\_t (triode PMOS) & 20 \% & −191 / −72 mV \\
unit\_w (well input) & 7 \% & +29 / +38 mV \\
unit\_q (undegenerated 2/6, square law, no R) & 22
\% & +75 / +200 mV \\
\end{longtable}

So the criterion for an ELIN unit here is not its
linearity but that its output depends on gp − gn
only. The split-tail pair meets it because each
device is a constant-current source follower
(Minch thesis ch. 7: a follower is linear whatever
the device law) and the element between the
sources sees exactly gp − gn. Any element whose
value depends on the absolute source voltage
(triode R, well κ, tail/CLM of a single-tail pair)
fails.

\hypertarget{1456--cascoding}{%
\section{14:56 ---
cascoding}\label{1456--cascoding}}

\textbf{(a) output devices = ESD clamp frames.} A
series cascode would put its drain, not the
clamp\textquotesingle s, on the pad; a
stacked-clamp frame would be a new ESD structure
to qualify. So no output cascode. What remains is
cascoding the internal high-impedance nodes.
Lengthening the folded and mirror cascodes (L 1 →
3 µm, W scaled, same overdrive) in
\texttt{d2s\_mp2}:

\begin{longtable}[]{@{}>{\raggedright\arraybackslash}p{(\linewidth - 10\tabcolsep) * \real{0.2493}}>{\raggedright\arraybackslash}p{(\linewidth - 10\tabcolsep) * \real{0.2332}}>{\raggedright\arraybackslash}p{(\linewidth - 10\tabcolsep) * \real{0.1763}}>{\raggedright\arraybackslash}p{(\linewidth - 10\tabcolsep) * \real{0.1649}}>{\raggedright\arraybackslash}p{(\linewidth - 10\tabcolsep) * \real{0.1763}}@{}}
\toprule\noalign{}
cascode L & T0 1 k ∥ 100 p & T0 100 p & T0 50 Ω &
PM 100 p \\
\midrule\noalign{}
\endhead
\bottomrule\noalign{}
\endlastfoot
1 µm (d2s\_mp2) & 69.7 dB & 99.6 dB & 43.9 dB &
63.4° \\
2 µm & 75.4 & 103.1 & 49.6 & 62.8° \\
3 µm (d2s\_mp2c3) & 78.9 & 104.3 & 53.1 & 62.0° \\
\end{longtable}

+9 dB for $\sim$320 µm² of gate area.
Regulated (gain-boosted) cascodes were not tried:
node a sits at $\sim$2.6 V and b at
$\sim$0.7 V (the output gates), which
leaves $\sim$0.35 V per device for the
mirror, the sinks and their cascodes, so no
headroom for an auxiliary
amplifier\textquotesingle s V\_GS.

\textbf{(b) separate output devices in the slot.}
Then they can be cascoded. Test case
\texttt{d2s\_lc2}: load-compensated (topology B,
no compensation capacitor at all), MP-DDA front
end, half-frame output devices, cascode gates
simply at vref (output range
$\sim$1.0\ldots2.4 V):

\begin{longtable}[]{@{}>{\raggedright\arraybackslash}p{(\linewidth - 16\tabcolsep) * \real{0.1444}}>{\raggedright\arraybackslash}p{(\linewidth - 16\tabcolsep) * \real{0.1208}}>{\raggedright\arraybackslash}p{(\linewidth - 16\tabcolsep) * \real{0.1292}}>{\raggedright\arraybackslash}p{(\linewidth - 16\tabcolsep) * \real{0.1119}}>{\raggedright\arraybackslash}p{(\linewidth - 16\tabcolsep) * \real{0.1119}}>{\raggedright\arraybackslash}p{(\linewidth - 16\tabcolsep) * \real{0.1119}}>{\raggedright\arraybackslash}p{(\linewidth - 16\tabcolsep) * \real{0.1582}}>{\raggedright\arraybackslash}p{(\linewidth - 16\tabcolsep) * \real{0.1119}}@{}}
\toprule\noalign{}
& T0 & 5 p & 20 p & 100 p & 1 n & 1 k ∥ 100 p &
I\_DD \\
\midrule\noalign{}
\endhead
\bottomrule\noalign{}
\endlastfoot
no cascode & 37.2 dB & 16° & 47° & 78° & 89° & no
loop gain & 174 µA \\
cascoded & 74.0 dB & unstable & 39° & 77° & 89° &
no loop gain & 169 µA \\
\end{longtable}

Cascoding adds 37 dB; a resistive load still
collapses a single-stage loop, and the minimum
load capacitance stays ($\sim$30--40 pF).
In the slot, the pad goes through IOPadAnalog,
whose \texttt{padres} is behind 587 Ω
(SecondaryProtection). Sensing at \texttt{padres}
(near side) makes that resistor an isolation
resistor: PM at 20 pF 39° → 62°, 100 pF → 104°;
the pad then sees a 587 Ω · C\_L low-pass outside
the loop (2.7 MHz at 100 pF). Sensing at
\texttt{pad} keeps the PM of the no-resistor case.
50 Ω through 587 Ω is not drivable at all. Open
point: the cascoded version has 7.8 mV systematic
offset (diode replicas at a different V\_DS than
the cascoded output devices); cascoding the diode
replicas would fix it.

\hypertarget{1458--noise-1-k--100-p-output-referred}{%
\section{14:58 --- noise (1 k ∥ 100 p,
output-referred)}\label{1458--noise-1-k--100-p-output-referred}}

\begin{longtable}[]{@{}>{\raggedright\arraybackslash}p{(\linewidth - 8\tabcolsep) * \real{0.3270}}>{\raggedright\arraybackslash}p{(\linewidth - 8\tabcolsep) * \real{0.2214}}>{\raggedright\arraybackslash}p{(\linewidth - 8\tabcolsep) * \real{0.2111}}>{\raggedright\arraybackslash}p{(\linewidth - 8\tabcolsep) * \real{0.2406}}@{}}
\toprule\noalign{}
& 1 kHz & 100 kHz & 10 Hz--10 MHz \\
\midrule\noalign{}
\endhead
\bottomrule\noalign{}
\endlastfoot
baseline A & 1599 nV/√Hz & 200 nV/√Hz & 278 µV
rms \\
MP-DDA (baseline sizing) & 1592 & 201 & 279 \\
d2s\_mp2c3 & 676 & 152 & 231 \\
\end{longtable}

\hypertarget{1459--d2s_mp2c3--d2s_mp2-with-l--3-uxb5m-cascodes-mos-caps-16--16-uxb5m-full-table}{%
\section{\texorpdfstring{14:59 ---
\texttt{d2s\_mp2c3} (= d2s\_mp2 with L = 3 µm
cascodes, MOS caps 16 × 16 µm): full
table}{14:59 --- d2s\_mp2c3 (= d2s\_mp2 with L = 3 µm cascodes, MOS caps 16 × 16 µm): full table}}\label{1459--d2s_mp2c3--d2s_mp2-with-l--3-uxb5m-cascodes-mos-caps-16--16-uxb5m-full-table}}

Renamed \texttt{d2s\_mpdda}
(\texttt{sim/d2s\_mpdda.spice}); every number
re-run in \texttt{sim/results\_improvements.txt}.
296 µA (baseline 427), T0 78.9 dB at 1 k ∥ 100 p
(66.5), 53.1 dB at 50 Ω (40.7), 104 dB capacitive
(96.5), PM 62° at 100 pF, 42° at 300 pF. Gain
0.4992\ldots0.5000 over 11 corners, nonlin ≤ 3.5
mV (ss 27 °C). Offset σ 5.3 mV, gain σ 1.1 \%
(rhigh area). THD 10 kHz −58.5 dB; 100 kHz −51.3
dB (1 k ∥ 100 p; lower fc than the MOM-MP version,
1.37 vs 1.62 MHz).

\hypertarget{15051520--consolidation}{%
\section{15:05--15:20 ---
consolidation}\label{15051520--consolidation}}

Intermediate netlists and scripts folded into one
runner, \texttt{sim/run\_improvements.py}, which
reuses \texttt{../run\_d2s.py} unchanged and
reproduces every number of this log and of
\texttt{proposed\_improvements.md} in one run (3
min 49 s, \texttt{sim/results\_improvements.txt}).
The re-run matches the numbers above. Files in
\texttt{sim/}:

\begin{longtable}[]{@{}>{\raggedright\arraybackslash}p{(\linewidth - 4\tabcolsep) * \real{0.3818}}>{\raggedright\arraybackslash}p{(\linewidth - 4\tabcolsep) * \real{0.6182}}@{}}
\toprule\noalign{}
file & content \\
\midrule\noalign{}
\endhead
\bottomrule\noalign{}
\endlastfoot
\texttt{units.spice} & DDA units: unit\_r (rhigh),
unit\_r2 (final, 2 × 5 µA), unit\_t (triode MOS),
unit\_w (well input), unit\_q (undegenerated) \\
\texttt{d2s\_mp.spice} & MP-DDA in the
baseline\textquotesingle s sizing; compensation
through the \texttt{ccomp} wrapper \\
\texttt{ccomp.spice} & ccomp\_mom (as
d2s\_miller), ccomp\_mos (hv PMOS accumulation
caps) \\
\texttt{d2s\_mpdda.spice} & the proposal (frames
kept as output devices) \\
\texttt{d2s\_bias\_lp.spice} & bias fixture with
26 µA of reference current instead of 90 µA \\
\texttt{d2s\_lc2.spice} & case (b): in-slot output
devices, load-compensated, cascoded (and not) \\
\texttt{run\_improvements.py},
\texttt{results\_improvements.txt} & runner and
its output \\
\end{longtable}

\hypertarget{1518--transfer-check}{%
\section{15:18 --- transfer
check}\label{1518--transfer-check}}

Checksums of the committed files against the
container copies:
\texttt{proposed\_improvements.md},
\texttt{sim/results\_improvements.txt} and
\texttt{sim/run\_improvements.py} did not match
after the first commits, although the tool
reported "written". The device still held earlier
versions of them, 20530 / 16367 / 13939 bytes
instead of 20571 / 18250 / 16300. Re-committing
from the same staged paths, with force, changed
nothing. Committing copies from a \emph{new}
staged directory (\texttt{outputs/xfer2/}) worked.
All files now match by sha256. This looks like the
silent no-op described in the project CLAUDE.md,
§6. Workaround: stage every new version of a file
under a new staged path.

\hypertarget{1520--time-and-tokens-used-by-this-task}{%
\section{15:20 --- time and tokens used by this
task}\label{1520--time-and-tokens-used-by-this-task}}

Wall clock: 14:08 → $\sim$15:20, about 72
minutes, interrupted twice by questions. Of that,
the compute (not the model) took about 3 min for
the OpenVAF build and model compile, about 6 min
for OCR on your computer, and about 10 min of
ngspice in total, including the 3 min 49 s full
re-run.

Tokens, from this session\textquotesingle s own
transcript
(\texttt{\textasciitilde{}/.claude/projects/-home-claude/\textless{}session\textgreater{}.jsonl},
per API response, as of 15:19):

\begin{longtable}[]{@{}>{\raggedright\arraybackslash}p{(\linewidth - 12\tabcolsep) * \real{0.3167}}>{\raggedright\arraybackslash}p{(\linewidth - 12\tabcolsep) * \real{0.1405}}>{\raggedright\arraybackslash}p{(\linewidth - 12\tabcolsep) * \real{0.1217}}>{\raggedright\arraybackslash}p{(\linewidth - 12\tabcolsep) * \real{0.1462}}>{\raggedright\arraybackslash}p{(\linewidth - 12\tabcolsep) * \real{0.1405}}>{\raggedright\arraybackslash}p{(\linewidth - 12\tabcolsep) * \real{0.1345}}@{}}
\toprule\noalign{}
phase & wall clock & API calls & output tokens &
cache writes & cache reads \\
\midrule\noalign{}
\endhead
\bottomrule\noalign{}
\endlastfoot
orientation: existing note, netlists, CACE
results, CLAUDE.md & 14:08--14:12 & 14 & 20,138 &
106,368 & 1,958,053 \\
simulation environment + literature reading/OCR,
interleaved & 14:12--14:27 & 49 & 23,126 & 122,848
& 12,499,286 \\
baseline reproduced; Caltech-thesis check &
14:27--14:31 & 6 & 2,887 & 13,684 & 1,960,595 \\
design reasoning + simulations & 14:31--15:00 & 40
& 125,771 & 157,464 & 17,134,286 \\
accounting, log timestamp correction &
15:00--15:02 & 5 & 10,264 & 17,604 & 2,485,393 \\
consolidated runner, re-run, notes, proposal,
transfer checks & 15:02--15:19 & 37 & 60,208 &
98,924 & 21,839,596 \\
\textbf{total, main thread} & & 151 & 242,394 &
517,090 & 56,877,712 \\
subagent (Caltech PDF check) & 14:30--14:31 & 5 &
1,463 & 94,882 & 368,538 \\
\end{longtable}

How to read it:

\begin{itemize}
\tightlist
\item
  \textbf{Output tokens} ($\sim$244 k):
  what the model wrote, including its reasoning.
  More than half of it falls in the design phase.
\item
  \textbf{Cache reads} ($\sim$57 M): each
  API call re-reads the whole conversation so far
  from the prompt cache. The context grew from
  $\sim$118 k tokens at the start (system
  prompt, tools, memory) to $\sim$608 k
  at the end. The growth came mostly from tool
  output: paper text, OCR, netlists and simulation
  listings.
\item
  \textbf{Cache writes} ($\sim$0.5 M):
  new context added to the cache.
\item
  \textbf{Uncached input}: 302 tokens in all.
\end{itemize}

How cache reads, cache writes and output count
against a specific plan\textquotesingle s usage
limits is not something this log can state (not
checked here); the plans\textquotesingle{} usage
limits are documented at support.claude.com. What
this kind of task pushes is the context size and
the number of calls made over it.

What drove the size, for planning similar tasks:

\begin{itemize}
\tightlist
\item
  Reading papers as text into the conversation:
  Edwards, Minch and Hasler excerpts and the
  Sarpeshkar OCR, roughly 100--150 k tokens of
  context.
\item
  Long simulation listings.
\item
  A long reasoning phase before the first
  simulation.
\end{itemize}

A second session that starts from
\texttt{proposed\_improvements.md} and
\texttt{log.md} instead of the papers would begin
with far less context.

\hypertarget{1555--ieeetran-version}{%
\section{15:55 --- IEEEtran
version}\label{1555--ieeetran-version}}

\texttt{proposed\_improvements.tex} is
\texttt{proposed\_improvements.md} in IEEEtran
journal style: the same content and numbers, with
equations and a references list taken from
\texttt{literature\_notes.md}.
\texttt{proposed\_improvements.pdf} (5 pages) was
built in the cloud with pdflatex, run twice, and
had no errors and no overfull boxes. The only
packages it needs besides IEEEtran.cls are
amsmath, booktabs, tabularx, array, textcomp, url
and cite; there is no siunitx. The sandbox I reach
on the linked computer has pdflatex but neither
IEEEtran.cls nor siunitx. Building there needs TeX
Live\textquotesingle s
\texttt{texlive-publishers}, or a copy of
IEEEtran.cls next to the file.

\hypertarget{1606--time-and-tokens-updated-through-the-latex-work}{%
\section{16:06 --- time and tokens, updated
through the LaTeX
work}\label{1606--time-and-tokens-updated-through-the-latex-work}}

Same method as at 15:20, now including the report,
the questions afterwards and the two LaTeX
conversions:

\begin{longtable}[]{@{}>{\raggedright\arraybackslash}p{(\linewidth - 12\tabcolsep) * \real{0.3075}}>{\raggedright\arraybackslash}p{(\linewidth - 12\tabcolsep) * \real{0.1424}}>{\raggedright\arraybackslash}p{(\linewidth - 12\tabcolsep) * \real{0.1233}}>{\raggedright\arraybackslash}p{(\linewidth - 12\tabcolsep) * \real{0.1482}}>{\raggedright\arraybackslash}p{(\linewidth - 12\tabcolsep) * \real{0.1424}}>{\raggedright\arraybackslash}p{(\linewidth - 12\tabcolsep) * \real{0.1363}}@{}}
\toprule\noalign{}
phase & wall clock & API calls & output tokens &
cache writes & cache reads \\
\midrule\noalign{}
\endhead
\bottomrule\noalign{}
\endlastfoot
orientation & 14:08--14:12 & 14 & 20,138 & 106,368
& 1,958,053 \\
simulation environment + literature/OCR &
14:12--14:27 & 49 & 23,126 & 122,848 &
12,499,286 \\
baseline reproduced; Caltech-thesis check &
14:27--14:31 & 6 & 2,887 & 13,684 & 1,960,595 \\
design reasoning + simulations & 14:31--15:00 & 40
& 125,771 & 157,464 & 17,134,286 \\
accounting, log timestamp correction &
15:00--15:02 & 5 & 10,264 & 17,604 & 2,485,393 \\
runner, re-run, notes, proposal, transfer checks,
report & 15:02--15:20 & 41 & 64,333 & 102,341 &
23,278,772 \\
questions about the spreadsheet button &
15:20--15:52 & 1 & 427 & 6,984 & 613,327 \\
IEEEtran LaTeX of the proposal & 15:52--16:04 & 16
& 32,152 & 45,403 & 10,226,268 \\
resource-log question; start of the
log\textquotesingle s LaTeX & 16:04--16:06 & 2 &
1,576 & 7,552 & 1,339,858 \\
\textbf{total, main thread} & & 174 & 280,674 &
580,248 & 71,495,838 \\
subagent (Caltech PDF check) & 14:30--14:31 & 5 &
1,463 & 94,882 & 368,538 \\
\end{longtable}

The context is now $\sim$674 k tokens per
call. The LaTeX conversion of this log comes after
this table and is not in it.

\hypertarget{1615--latex-version-of-this-log}{%
\section{16:15 --- LaTeX version of this
log}\label{1615--latex-version-of-this-log}}

\texttt{log.tex} / \texttt{log.pdf} (article
class, not IEEEtran, in the style of
\texttt{sudelbuecher/esd\_protection/*.tex}). The
body was converted from this file with pandoc
3.1.3, and the tables were given widths weighted
by their content. Builds with pdflatex using only
standard packages (booktabs, longtable, calc,
hyperref, xurl, microtype, textcomp, geometry). A
snapshot: entries after 16:15 are only in
\texttt{log.md}.

\hypertarget{1637--testbench-decks-in-simtb}{%
\section{\texorpdfstring{16:37 --- testbench decks
in
\texttt{sim/tb/}}{16:37 --- testbench decks in sim/tb/}}\label{1637--testbench-decks-in-simtb}}

Nine stand-alone ngspice decks, so the results can
be rerun without \texttt{run\_improvements.py}:
\texttt{tb\_mpdda\_dc}, \texttt{tb\_mpdda\_loop},
\texttt{tb\_mpdda\_step}, \texttt{tb\_mpdda\_thd},
\texttt{tb\_mpdda\_noise}, \texttt{tb\_mpdda\_op},
\texttt{tb\_units}, \texttt{tb\_moscv} and
\texttt{tb\_lc2\_loop}. Each one names its
reference value from
\texttt{results\_improvements.txt} in the header
and runs from its own directory
(\texttt{ngspice\ -b\ \textless{}deck\textgreater{}}
prints,
\texttt{ngspice\ \textless{}deck\textgreater{}}
plots). They need the PDK \texttt{.spiceinit}
(sourcepath to the models, OSDI for psp103,
r3\_cmc and cap\_cmomi). Rerun in the cloud, they
reproduce the reference: gain 0.49994, offset
0.055 mV, nonlinearity 1.243 mV, Idd 296 µA; loop
78.9 dB, 1.366 MHz, 72.4°; THD 0.119 \%; noise 676
/ 152 nV/√Hz at 1 / 100 kHz; \texttt{d2s\_lc2}
near side 74.0 dB, 1.49 MHz, 104.3°.

\hypertarget{1700--xschem-schematics-in-xschem}{%
\section{\texorpdfstring{17:00 --- xschem
schematics in
\texttt{xschem/}}{17:00 --- xschem schematics in xschem/}}\label{1700--xschem-schematics-in-xschem}}

First drafts for hand editing, laid out like the
hand-edited \texttt{../xschem/d2s\_miller.sch}:
ports in one column on the left, output on the
right, vdd and vss as wires across the sheet, one
column per current branch, bulks wired, and the
input and bias nets as wire buses from the port
column. Internal nets carry \texttt{lab\_wire}
labels so that the netlist keeps the
source\textquotesingle s net names.

\begin{itemize}
\tightlist
\item
  \texttt{unit\_r2} (the unit of
  \texttt{d2s\_mpdda}), and \texttt{unit\_r},
  \texttt{unit\_t}, \texttt{unit\_w},
  \texttt{unit\_q}. All five share one box symbol:
  gp and gn on the left, vdd and vbp on top, x, y
  and vss at the bottom.
\item
  \texttt{d2s\_mpdda}: four \texttt{unit\_r2}
  boxes; the fold, mirror, class-AB pair and
  output devices sit where they are in
  \texttt{d2s\_miller.sch}. CMA and CMB are drawn
  as hv PMOS accumulation capacitors.
\item
  \texttt{d2s\_lc2} and \texttt{d2s\_lc2\_nc}: the
  same core. The output cascode gates reach vref
  through one \texttt{lab\_pin}.
\item
  \texttt{d2s\_bias\_lp}, and
  \texttt{d2s\_mpdda\_biased}, the CACE fixture,
  which is wired like
  \texttt{d2s\_miller\_biased.sch}.
\item
  The symbols of \texttt{d2s\_mpdda},
  \texttt{d2s\_lc2(\_nc)}, \texttt{d2s\_bias\_lp}
  and \texttt{d2s\_mpdda\_biased} are copies of
  \texttt{d2s\_miller.sym}, \texttt{d2s\_bias.sym}
  and \texttt{d2s\_miller\_biased.sym}, which have
  the same port lists.
\end{itemize}

Parameters are frozen at the \texttt{.subckt}
defaults:

\begin{itemize}
\tightlist
\item
  lcas = 3, so CX, CY, PCL and PCR are 30u / 3u
  with ng = 2.
\item
  wc = lc = 16u.
\item
  rl = 50.6u in \texttt{unit\_r2} and 36.4u in
  \texttt{unit\_r}.
\item
  iab = 5u, ibnc = 2u, wdp = 6.66u, wdn = 3.8u.
\item
  \texttt{unit\_w} keeps \texttt{w=wwi\ l=lwi},
  which the testbench defines.
\end{itemize}

Not drawn: \texttt{d2s\_mp}, whose \texttt{ccomp}
wrapper exists only in the testbench, and
\texttt{ccomp\_mom} / \texttt{ccomp\_mos};
\texttt{ccomp\_mos} appears in \texttt{d2s\_mpdda}
as CMA and CMB.

\texttt{check\_xschem.py} netlists every
\texttt{.sch} with xschem and compares each
subcircuit in the netlist with
\texttt{../sim/*.spice}, device by device. It
checks model, nets in order, w / l / ng / value,
and the port order of both \texttt{.sch} and
\texttt{.sym}. With xschem 3.4.4 it finds 0
mismatches over the 10 cells, including the
\texttt{unit\_r2} bodies inside the parents. It
did report each deliberate error I introduced: a
changed port label, a resistor length, a current
value and two swapped \texttt{.sym} pins.

In the first draft, nets l2 and b came out split
(net1, net2). A bus crossed the middle of a column
without a junction, and xschem connects a wire
only at its end points. Hand edits can break
connectivity the same way, so rerun
\texttt{check\_xschem.py} after editing. Most of
the 106 wire crossings in \texttt{d2s\_mpdda} are
in the input and bias bus band under the four
units; that band is the obvious place for hand
rearrangement or for labels. The script that drew
the sheets stays in the cloud, since the drafts
are meant to be edited by hand from here on.

\hypertarget{1715--latex-log-brought-up-to-date}{%
\section{17:15 --- LaTeX log brought up to
date}\label{1715--latex-log-brought-up-to-date}}

\texttt{log.tex} and \texttt{log.pdf} were
regenerated from this file through this entry,
with the same steps as at 16:15 (pandoc, then the
table widths). They are still snapshots: entries
after this one are only in \texttt{log.md} until
the next regeneration.

\hypertarget{1735--block-annotations-flat-schematics-xschemreadmemd}{%
\section{\texorpdfstring{17:35 --- block
annotations, flat schematics,
\texttt{xschem/README.md}}{17:35 --- block annotations, flat schematics, xschem/README.md}}\label{1735--block-annotations-flat-schematics-xschemreadmemd}}

Every sheet now outlines its functional blocks
with magenta dashed boxes tagged {[}1{]} to
{[}9{]}, with a legend under the sheet. A number
means the same block on every sheet: {[}1{]} DDA
units, {[}2{]} fold, {[}3{]} mirror, {[}4{]}
class-AB control, {[}5{]} its copy in the mirror
input branch, {[}6{]} output devices, {[}7{]}
Miller capacitors (cascodes in \texttt{d2s\_lc2}),
{[}8{]} bias, {[}9{]} diode replicas in
\texttt{d2s\_lc2}. \texttt{xschem/README.md}
explains the blocks with sizes, nets and tt
currents, gives the signal path and the bias
network, and lists the files.

Two new sheets show everything at transistor
level:

\begin{itemize}
\tightlist
\item
  \texttt{d2s\_mpdda\_flat}: \texttt{d2s\_mpdda}
  with the four units drawn out. Unit devices and
  nets carry the unit\textquotesingle s name
  (\texttt{Ta\_A}, \texttt{sa\_A}, ...). It has
  the same 13 pins.
\item
  \texttt{d2s\_mpdda\_bias\_flat}:
  \texttt{d2s\_mpdda} and \texttt{d2s\_bias\_lp}
  on one sheet, with the bias columns between the
  port column and the units. It has 7 pins; the
  bias nets are internal.
\end{itemize}

\texttt{check\_xschem.py} compares the flat sheets
with the sources with the subcircuits expanded
under that naming. 0 mismatches over 12 cells. A
swapped unit-net label and a changed rhigh body
were both reported.

\hypertarget{1851--testbench-schematics-and-figures}{%
\section{18:51 --- testbench schematics and
figures}\label{1851--testbench-schematics-and-figures}}

\textbf{Testbench schematics.}
\texttt{xschem/tb\_*.sch} has one sheet per deck
in \texttt{sim/tb/}, nine in all. They are laid
out like \texttt{../xschem/tb\_d2s\_*.sch}. The
DUT and the bias come from the schematics in
\texttt{xschem/}. Each sheet\textquotesingle s
code block holds the deck\textquotesingle s
\texttt{.lib}, \texttt{.param} and \texttt{.save}
lines and its \texttt{.control} section,
unchanged. \texttt{check\_xschem.py} now also
compares every testbench with its deck: each
element (nets with GND = 0, value or model and
sizes) and every code line.

Result: 0 mismatches over 12 cells and 9
testbenches. A changed RL value and a changed
\texttt{meas} line were both reported. Netlisted
by xschem and run in ngspice, all nine reproduce
the decks\textquotesingle{} reference values: loop
78.9 dB / 1.366 MHz / 72.4°, gain 0.49994, THD
0.1187 \%, noise 230 µV, units 5.442 / 5.416 µA,
MOS C--V 0.370 / 0.648 / 0.745 pF, and
\texttt{d2s\_lc2} 74.0 dB / 1.49 MHz / 104.3°.

\textbf{Sheets re-saved between 18:00 and 18:16.}
Seven sheets on the linked computer had been saved
again with xschem 3.4.8RC: \texttt{d2s\_bias\_lp},
\texttt{d2s\_lc2}, \texttt{d2s\_lc2\_nc},
\texttt{d2s\_mpdda}, \texttt{d2s\_mpdda\_flat},
\texttt{d2s\_mpdda\_bias\_flat} and
\texttt{unit\_r2}. The differences are
xschem\textquotesingle s own rewriting: merged
wire segments, reversed wire directions and
reordered texts. Wire geometry and instances are
unchanged, and the checker passes on these
versions. They were left as they are; the figures
were made from them.

\textbf{Figures.}
\texttt{figures/\textless{}sheet\textgreater{}.svg},
\texttt{.pdf} and \texttt{.png} exist for all 21
sheets (63 files). Each is
xschem\textquotesingle s own SVG export with a
white background and black wires, symbols and
text. Block frames and legends are \#d55e00, about
47 \% grey in monochrome. Each figure is cropped
to the drawing and leaves out the title block, pin
letters, pin squares, \texttt{m=1}, \texttt{ng=1},
\texttt{b=0}, the rhigh \texttt{R=} expression and
the MOS model names. The PNGs are 2 px per unit,
at most 6000 px wide; the PDFs embed Liberation
Sans. Details are in \texttt{xschem/README.md}.

The SVG and PNG files reached the linked computer
with a C2PA provenance block added in transfer:
\texttt{\textless{}metadata\textgreater{}} in the
SVG and a \texttt{caBX} chunk in the PNG. The
pictures themselves are unchanged (pixel
difference 0), and the PDFs arrived unchanged.

\hypertarget{1859--regeneration-scripts-in-xschemscripts}{%
\section{\texorpdfstring{18:59 --- regeneration
scripts in
\texttt{xschem/scripts/}}{18:59 --- regeneration scripts in xschem/scripts/}}\label{1859--regeneration-scripts-in-xschemscripts}}

\begin{itemize}
\tightlist
\item
  \texttt{export\_figures.py} makes the figures;
  by default it does all sheets of
  \texttt{xschem/} into \texttt{figures/}. It
  needs xschem, \texttt{\$PDK\_ROOT}, cairosvg and
  Pillow.
\item
  \texttt{gen\_cells.py\ OUTDIR} and
  \texttt{gen\_testbenches.py\ OUTDIR} drew the
  cell and testbench sheets; they share the helper
  \texttt{xsheet.py}. They refuse to write into
  \texttt{xschem/} unless \texttt{-\/-overwrite}
  is given, because the sheets there are now
  edited by hand.
\end{itemize}

Before shipping, all three ran in the cloud from a
copy of the folder layout. The figures came out
identical to those in \texttt{figures/}: the same
SVG, and PNG pixels with no difference. The
generators\textquotesingle{} 34 output files were
byte-identical to the first drafts. The guard
against writing into \texttt{xschem/} also worked.
The scripts have not been run on the linked
computer; its sandbox has no xschem.
\texttt{xschem/README.md} has a "Regenerating"
section with the commands.

\end{document}
